% Pregunta 814 (2005-septiembre-OEP2005, n.º 33): arancel en un país pequeño, equilibrio parcial.
% El precio interior pasa de P* a P*+t: la producción sube (S0 -> S1), el consumo baja (D0 -> D1),
% las importaciones se reducen (M0 -> M1), el Estado recauda t*M1 y quedan dos triángulos de pérdida.
% Curvas: S: P = 0,5 + 0,8Q; D: P = 4,6 - 0,7Q; P* = 1,3; P*+t = 2,0.
% Autor: Víctor Gutiérrez Marcos
\documentclass[tikz,border=6pt]{standalone}
\input{tikz-tcee}
\begin{document}
\begin{tikzpicture}[tcee]
  \ejes{$Q$}{$P$}{6.6}{5}
  % Áreas
  \fill[area] (1.875,1.3) rectangle (3.714,2.0);
  \fill[area2] (1.0,1.3) -- (1.875,1.3) -- (1.875,2.0) -- cycle;
  \fill[area2] (3.714,1.3) -- (4.714,1.3) -- (3.714,2.0) -- cycle;
  \node[etiqueta,text=tceemorado,font=\footnotesize] at (2.795,1.65) {recaudación};
  % Precios
  \draw[guia,thick] (0,1.3) -- (6.2,1.3);
  \draw[curvanueva] (0,2.0) -- (6.2,2.0);
  \node[etiqueta,left] at (0,1.3) {$P^*$};
  \node[etiqueta,left,text=tceerojo] at (0,2.0) {$P^*\!+t$};
  % Curvas
  \draw[curva1] (0,0.5) -- (5.0,4.5) node[etiqueta,above] {$S$};
  \draw[curva2] (0,4.6) -- (6.0,0.4) node[etiqueta,right] {$D$};
  % Proyecciones
  \draw[guia] (1.0,1.3) -- (1.0,0) node[below,text=black] {$S_0$};
  \draw[guia] (4.714,1.3) -- (4.714,0) node[below,text=black] {$D_0$};
  \draw[guia,tceerojo] (1.875,2.0) -- (1.875,0) node[below,text=tceerojo] {$S_1$};
  \draw[guia,tceerojo] (3.714,2.0) -- (3.714,0) node[below,text=tceerojo] {$D_1$};
  \node[punto] at (1.0,1.3) {};
  \node[punto] at (4.714,1.3) {};
  \node[punto,fill=tceerojo] at (1.875,2.0) {};
  \node[punto,fill=tceerojo] at (3.714,2.0) {};
  % Importaciones
  \draw[<->,thick,tceegris] (1.0,-0.75) -- node[fill=white,inner sep=1pt,font=\footnotesize] {$M_0$ (antes)} (4.714,-0.75);
  \draw[<->,thick,tceerojo] (1.875,-1.3) -- (3.714,-1.3);
  \node[etiqueta,below,font=\footnotesize,text=tceerojo] at (2.795,-1.3) {$M_1$ (con arancel)};
  % Etiquetas de los triángulos
  \node[etiqueta,align=center,font=\footnotesize,text=tceedorado!80!black] (pp) at (0.95,3.4) {pérdida de\\producción};
  \draw[->,thin,tceedorado!80!black] (pp.south) -- (1.6,1.55);
  \node[etiqueta,align=center,font=\footnotesize,text=tceedorado!80!black] (pc) at (5.4,2.75) {pérdida de\\consumo};
  \draw[->,thin,tceedorado!80!black] (pc.south) -- (4.05,1.55);
\end{tikzpicture}
\end{document}
